% ============================================================
% 测试卷 — 第2章 一元二次函数、方程与不等式
% ============================================================
\documentclass[12pt, a4paper]{ctexart}
\usepackage[margin=2cm]{geometry}
\usepackage{amsmath, amssymb}
\usepackage{enumitem}
\usepackage{fancyhdr}
\usepackage{xcolor}
\usepackage{tikz}
\usepackage{graphicx}
\usepackage{tabularx}
\usepackage{multicol}

\definecolor{mainblue}{RGB}{41, 98, 168}

\pagestyle{fancy}
\fancyhf{}
\fancyhead[L]{\small\textcolor{mainblue}{2025-2026学年度·第2章测试卷}}
\fancyhead[R]{\small\textcolor{mainblue}{数学（必修第一册）}}
\fancyfoot[C]{\small —\ \thepage\ —}
\renewcommand{\headrulewidth}{0.6pt}
\renewcommand{\headrule}{\hbox to\headwidth{\color{mainblue}\leaders\hrule height \headrulewidth\hfill}}

\setCJKmainfont{Hiragino Sans GB}
\setCJKsansfont{Heiti SC}

\newcommand{\blank}[1]{\underline{\hspace{#1}}}

\begin{document}

\begin{center}
{\Large\bfseries 第2章\quad 一元二次函数、方程与不等式\quad 测试卷}\\[6pt]
{\normalsize 数学（必修第一册）}\\[4pt]
{\small 考试时间：120分钟\quad 满分：150分}
\end{center}

\vspace{2pt}
\begin{center}
\begin{tabular}{|c|c|c|c|c|}
\hline
题号 & 一 & 二 & 三 & 总分 \\ \hline
得分 & & & & \\ \hline
\end{tabular}
\end{center}
\vspace{4pt}

\noindent\textbf{注意事项：}
\begin{enumerate}[leftmargin=2em]
  \item 答题前，请在答题卡上填写姓名、班级和考号。
  \item 选择题请用2B铅笔在答题卡上涂黑对应选项。
  \item 非选择题请用黑色签字笔在答题卡指定区域作答。
\end{enumerate}

\vspace{6pt}

% ============ 一、单项选择题 ============
\noindent\rule{\textwidth}{0.4pt}\\[2pt]
{\large\bfseries 一、单项选择题}（本大题共8小题，每小题5分，共40分。每小题只有一个选项符合题目要求。）

\vspace{6pt}

\begin{enumerate}
\item 方程 $x^2-5x+6=0$ 的解集为（\quad）
  \begin{enumerate}[label=\Alph*., itemsep=0pt]
    \item $\{2\}$ \quad \item $\{3\}$ \quad \item $\{2,3\}$ \quad \item $\{-2,-3\}$
  \end{enumerate}

\item 一元二次方程 $x^2+2x+5=0$ 的判别式 $\Delta$ 的值为（\quad）
  \begin{enumerate}[label=\Alph*., itemsep=0pt]
    \item $-16$ \quad \item $16$ \quad \item $-24$ \quad \item $24$
  \end{enumerate}

\item 不等式 $x^2-4x+3<0$ 的解集为（\quad）
  \begin{enumerate}[label=\Alph*., itemsep=0pt]
    \item $\{x\mid x<1 \text{ 或 } x>3\}$ \quad \item $\{x\mid 1<x<3\}$
    \item $\{x\mid x<1\}$ \quad \item $\{x\mid x>3\}$
  \end{enumerate}

\item 二次函数 $y=x^2-2x+3$ 的顶点坐标为（\quad）
  \begin{enumerate}[label=\Alph*., itemsep=0pt]
    \item $(1,2)$ \quad \item $(-1,2)$ \quad \item $(1,4)$ \quad \item $(-1,4)$
  \end{enumerate}

\item 若关于 $x$ 的方程 $x^2+2x+k=0$ 有两个不相等的实数根，则 $k$ 的取值范围为（\quad）
  \begin{enumerate}[label=\Alph*., itemsep=0pt]
    \item $k<1$ \quad \item $k\leq 1$ \quad \item $k>1$ \quad \item $k\geq 1$
  \end{enumerate}

\item 若 $x_1$、$x_2$ 是方程 $x^2-3x+1=0$ 的两个根，则 $x_1^2+x_2^2$ 的值为（\quad）
  \begin{enumerate}[label=\Alph*., itemsep=0pt]
    \item $9$ \quad \item $7$ \quad \item $5$ \quad \item $3$
  \end{enumerate}

\item 不等式 $-x^2+2x+3\geq 0$ 的解集为（\quad）
  \begin{enumerate}[label=\Alph*., itemsep=0pt]
    \item $\{x\mid -1\leq x\leq 3\}$ \quad \item $\{x\mid x\leq -1 \text{ 或 } x\geq 3\}$
    \item $\{x\mid -3\leq x\leq 1\}$ \quad \item $\{x\mid x\leq -3 \text{ 或 } x\geq 1\}$
  \end{enumerate}

\item 已知函数 $f(x)=ax^2+bx+c$（$a\neq 0$）的图象过点 $(0,1)$ 和 $(2,1)$，且在 $(-\infty,1)$ 上单调递减，则下列结论正确的是（\quad）
  \begin{enumerate}[label=\Alph*., itemsep=0pt]
    \item $a>0$，$b<0$ \quad \item $a>0$，$b>0$
    \item $a<0$，$b>0$ \quad \item $a<0$，$b<0$
  \end{enumerate}
\end{enumerate}

\vspace{6pt}

% ============ 二、多项选择题 ============
\noindent\rule{\textwidth}{0.4pt}\\[2pt]
{\large\bfseries 二、多项选择题}（本大题共4小题，每小题5分，共20分。全部选对得5分，部分选对得2分，有选错的得0分。）

\vspace{6pt}

\begin{enumerate}[resume]
\item 下列关于二次函数 $y=-x^2+4x-3$ 的说法正确的是（\quad）
  \begin{enumerate}[label=\Alph*., itemsep=0pt]
    \item 函数图象开口向下
    \item 函数的最大值为 $1$
    \item 对称轴为 $x=2$
    \item 函数在 $(2,+\infty)$ 上单调递增
  \end{enumerate}

\item 关于 $x$ 的方程 $x^2-(m+1)x+m=0$，下列结论正确的是（\quad）
  \begin{enumerate}[label=\Alph*., itemsep=0pt]
    \item 方程必有实数根
    \item 方程必有一根为 $1$
    \item 当 $m=1$ 时，方程有两个相等的实数根
    \item 当 $m=-1$ 时，方程的两根之积为 $1$
  \end{enumerate}

\item 已知 $a<b$，则下列不等式正确的是（\quad）
  \begin{enumerate}[label=\Alph*., itemsep=0pt]
    \item $a^2<b^2$
    \item $a+1<b+1$
    \item $-2a>-2b$
    \item $\dfrac{a}{3}<\dfrac{b}{3}$
  \end{enumerate}

\item 已知函数 $f(x)=x^2-2ax+1$ 在区间 $[1,+\infty)$ 上单调递增，则（\quad）
  \begin{enumerate}[label=\Alph*., itemsep=0pt]
    \item $a\leq 1$
    \item $f(0)>f(1)$ 不一定成立
    \item $f(a)$ 是函数的最小值
    \item 当 $a=1$ 时，$f(x)$ 在 $[1,+\infty)$ 上的最小值为 $0$
  \end{enumerate}
\end{enumerate}

\vspace{6pt}

% ============ 三、填空题 ============
\noindent\rule{\textwidth}{0.4pt}\\[2pt]
{\large\bfseries 三、填空题}（本大题共4小题，每小题5分，共20分。）

\vspace{6pt}

\begin{enumerate}[resume]
\item 方程 $2x^2-7x+3=0$ 的两根之和为 \underline{\hspace{3cm}}。

\item 不等式 $\dfrac{x-1}{x+2}\geq 0$ 的解集为 \underline{\hspace{4cm}}。

\item 若关于 $x$ 的不等式 $x^2+ax+b<0$ 的解集为 $\{x\mid 1<x<3\}$，则 $a+b=$ \underline{\hspace{2cm}}。

\item 已知函数 $f(x)=x^2-4x+5$，$x\in[0,3]$，则 $f(x)$ 的值域为 \underline{\hspace{4cm}}。
\end{enumerate}

\vspace{6pt}

% ============ 四、解答题 ============
\noindent\rule{\textwidth}{0.4pt}\\[2pt]
{\large\bfseries 四、解答题}（本大题共6小题，共70分。解答应写出文字说明、证明过程或演算步骤。）

\vspace{6pt}

\begin{enumerate}[resume]
\item[17.]（本小题10分）\\
解下列方程与不等式：
\begin{enumerate}[(1)]
  \item 解方程 $3x^2-10x+3=0$；
  \item 解不等式 $2x^2+x-6>0$。
\end{enumerate}

\vspace{6cm}

\item[18.]（本小题12分）\\
已知关于 $x$ 的方程 $x^2-2mx+m+2=0$。
\begin{enumerate}[(1)]
  \item 当 $m$ 取何值时，方程有两个不相等的实数根？
  \item 当 $m=3$ 时，求方程的两根 $x_1$、$x_2$，并求 $x_1^2+x_2^2$ 的值。
\end{enumerate}

\vspace{6cm}

\item[19.]（本小题12分）\\
已知二次函数 $f(x)=ax^2+bx+c$ 满足 $f(0)=1$，$f(1)=0$，且函数的最小值为 $-\dfrac{1}{4}$。
\begin{enumerate}[(1)]
  \item 求 $a$、$b$、$c$ 的值；
  \item 求 $f(x)$ 在区间 $[-1,2]$ 上的最大值和最小值。
\end{enumerate}

\vspace{6cm}

\item[20.]（本小题12分）\\
某商店以每件 $40$ 元的价格购进一种商品，若以每件 $50$ 元出售，每天可售出 $200$ 件。经调查发现，每件商品的售价每上涨 $1$ 元，每天的销售量就减少 $10$ 件。设每件商品的售价为 $x$ 元（$x>50$），每天的利润为 $y$ 元。
\begin{enumerate}[(1)]
  \item 写出 $y$ 关于 $x$ 的函数关系式；
  \item 当每件商品的售价定为多少元时，每天的利润最大？最大利润是多少？
\end{enumerate}

\vspace{6cm}

\item[21.]（本小题12分）\\
已知关于 $x$ 的不等式 $ax^2+bx+2>0$ 的解集为 $\{x\mid -1<x<2\}$。
\begin{enumerate}[(1)]
  \item 求 $a$、$b$ 的值；
  \item 解关于 $x$ 的不等式 $ax^2-bx+2>0$。
\end{enumerate}

\vspace{7cm}

\item[22.]（本小题12分）\\
已知函数 $f(x)=x^2-2ax+2$（$a\in\mathbb{R}$）。
\begin{enumerate}[(1)]
  \item 当 $a=1$ 时，求 $f(x)$ 在区间 $[-1,3]$ 上的值域；
  \item 若 $f(x)\geq a$ 对一切 $x\in\mathbb{R}$ 恒成立，求实数 $a$ 的取值范围。
\end{enumerate}

\vspace{7cm}
\end{enumerate}

\vfill
\begin{center}
\small\textcolor{gray}{本试卷仅供教学参考}
\end{center}

\end{document}
